\uuid{T9TC}
\titre{Changements de variables/éléments simples}
\auteur{hanani ; mzari}
\organisation{exo7}
\licence{CC BY-NC-SA 4.0}
\datecreate{2020-01-20}
\niveau{L1}
\module{Analyse}
\chapitre{Calcul d'intégrales}
\sousChapitre{Primitives diverses}
\difficulte{3}
\theme{Primitives, changements de variables, éléments simples}
\exo7qcmid{lille-2.primitives.048}

\contenu{
\begin{question}
\qclassification{127.11}
\qtags{motcle=changements de variables/éléments simples}
Parmi les affirmations suivantes, cocher celles qui sont vraies :
\begin{answers}
\bad{Le changement de variable $\displaystyle t=\sqrt{x}$ donne $\displaystyle \int \frac{\mathrm{d}x}{\sqrt{x}(x+1)^2}=\int \frac{\mathrm{d}t}{(t^2+1)^2}$.}
\good{$\displaystyle \int \frac{\mathrm{d}t}{t^2+1}=\frac{t}{t^2+1}+\int \frac{2t^2\, \mathrm{d}t}{(t^2+1)^2}$.}
\good{$\displaystyle \int \frac{2\, \mathrm{d}t}{(t^2+1)^2}=\frac{t}{t^2+1}+\arctan t+k$, $k\in \Rr$.}
\bad{Une primitive de $\displaystyle \frac{1}{\sqrt{x}(x+1)^2}$ sur $]0,+\infty[$ est $\displaystyle \frac{1}{2}\left(\frac{\sqrt{x}}{x+1}+\arctan \sqrt{x}\right)$.}
\end{answers}
\vskip2mm
\begin{explanations}
Avec $\displaystyle t=\sqrt{x}$, on a : $\displaystyle x=t^2\Rightarrow \mathrm{d}x=2t\, \mathrm{d}t$ et $\displaystyle \int \frac{\mathrm{d}x}{\sqrt{x}(x+1)^2}=\int \frac{2\, \mathrm{d}t}{(t^2+1)^2}$.
Une intégration par parties avec $\displaystyle u=\frac{1}{t^2+1}$ et $v=t$ donne
$$\int \frac{\mathrm{d}t}{t^2+1}=\frac{t}{t^2+1}+\int \frac{2t^2\, \mathrm{d}t}{(t^2+1)^2}=\frac{t}{t^2+1}+\int \frac{2\, \mathrm{d}t}{t^2+1}-\int \frac{2\, \mathrm{d}t}{(t^2+1)^2}.$$
D'où $\displaystyle \int \frac{2\, \mathrm{d}t}{(t^2+1)^2}=\frac{t}{t^2+1}+\arctan t+k\Rightarrow \int \frac{\mathrm{d}x}{\sqrt{x}(x+1)^2}=\frac{\sqrt{x}}{x+1}+\arctan \sqrt{x}+k$, $k\in \Rr$.
\end{explanations}
\end{question}
}
